% GENERATED FILE. Edit canonical lessons, then run npm run generate:books.
% canonical-source-hash: fnv1a64:00c9f82c711c1eda
% canonical-lessons: AR-C69-asir, AR-C69-dajaj, AR-C69-manzil, AR-C69-ghurfa, AR-C69-nafidha

\chapter{Juice, Chicken, House, Room, Window}
\label{ch:ar-juice-chicken-house-room-window}
\hlchaptermodality{69}

\begin{chapteropening}
\textbf{By the end of this chapter:} \emph{I can say juice, chicken, a house, a room and a window.}

Complete the last lesson of chapter 69: I can say juice, chicken, a house, a room and a window.
\end{chapteropening}

\section[\texorpdfstring{\ar{عصير} (ʿaṣīr)}{عصير (ʿaṣīr)}]{\ar{عصير} (ʿaṣīr) — juice}

\label{lesson:AR-C69-asir}



\begin{quote}
Before the new one: say the Arabic for onions, then the Arabic for breakfast.
\end{quote}

\subsection*{You'll want to know: \ar{عصير}}
\textbf{\ar{عصير}} — \emph{ʿaṣīr} — \textquotedblleft{}juice\textquotedblright{}.

\begin{cousinweb}
\textbf{\ar{عصير}} (\emph{ʿaṣīr}) is built on the native Arabic root \emph{ʿ-ṣ-r}, \textquotedblleft{}to press, to squeeze\textquotedblright{} (Lane). \emph{ʿAṣīr} is what is pressed out: juice.
\end{cousinweb}

\begin{grammarlens}[title={What you've built}]
One more word for this chapter.
\end{grammarlens}

\subsection*{Your turn}
\begin{itemize}
\raggedright
  \item \emph{Say it:} \emph{ʿaṣīr}
  \item \emph{Say it:} \emph{ʿaṣīr}, once more
  \item \emph{Say it:} say \emph{ʿaṣīr}
  \item \emph{Recall:} say the Arabic for onions, then the Arabic for breakfast, then say \emph{ʿaṣīr} again
\end{itemize}

\begin{tcolorbox}[breakable,colback=teal!4,colframe=teal!35!black,title={Before you move on}]
What does \ar{عصير} mean? (\textquotedblleft{}Juice\textquotedblright{}.) Say it once more.
\end{tcolorbox}
\section[\texorpdfstring{\ar{دجاج} (dajāj)}{دجاج (dajāj)}]{\ar{دجاج} (dajāj) — chicken}

\label{lesson:AR-C69-dajaj}



\begin{quote}
Before the new one: say the Arabic for breakfast, then the Arabic for juice.
\end{quote}

\subsection*{You'll want to know: \ar{دجاج}}
\textbf{\ar{دجاج}} — \emph{dajāj} — \textquotedblleft{}chicken\textquotedblright{}.

\begin{cousinweb}
\textbf{\ar{دجاج}} (\emph{dajāj}) is built on the native Arabic root \emph{d-j-j}, \textquotedblleft{}the domestic fowl\textquotedblright{} (Lane). Lane lists \emph{dajāj}, the common fowl, under \emph{d-j-j}.
\end{cousinweb}

\begin{grammarlens}[title={What you've built}]
One more word for this chapter.
\end{grammarlens}

\subsection*{Your turn}
\begin{itemize}
\raggedright
  \item \emph{Say it:} \emph{dajāj}
  \item \emph{Say it:} \emph{dajāj}, once more
  \item \emph{Say it:} say \emph{dajāj}
  \item \emph{Recall:} say the Arabic for breakfast, then the Arabic for juice, then say \emph{dajāj} again
\end{itemize}

\begin{tcolorbox}[breakable,colback=teal!4,colframe=teal!35!black,title={Before you move on}]
What does \ar{دجاج} mean? (\textquotedblleft{}Chicken\textquotedblright{}.) Say it once more.
\end{tcolorbox}
\section[\texorpdfstring{\ar{منزل} (manzil)}{منزل (manzil)}]{\ar{منزل} (manzil) — a house}

\label{lesson:AR-C69-manzil}



\begin{quote}
Before the new one: say the Arabic for juice, then the Arabic for chicken.
\end{quote}

\subsection*{You'll want to know: \ar{منزل}}
\textbf{\ar{منزل}} — \emph{manzil} — \textquotedblleft{}a house\textquotedblright{}.

\begin{cousinweb}
\textbf{\ar{منزل}} (\emph{manzil}) is built on the native Arabic root \emph{n-z-l}, \textquotedblleft{}to come down, to alight\textquotedblright{} (Lane). \emph{Manzil} is the \emph{ma-} place noun, \textquotedblleft{}a place of alighting, a lodging\textquotedblright{}.
\end{cousinweb}

\begin{grammarlens}[title={What you've built}]
One more word for this chapter.
\end{grammarlens}

\subsection*{Your turn}
\begin{itemize}
\raggedright
  \item \emph{Say it:} \emph{manzil}
  \item \emph{Say it:} \emph{manzil}, once more
  \item \emph{Say it:} say \emph{manzil}
  \item \emph{Recall:} say the Arabic for juice, then the Arabic for chicken, then say \emph{manzil} again
\end{itemize}

\begin{tcolorbox}[breakable,colback=teal!4,colframe=teal!35!black,title={Before you move on}]
What does \ar{منزل} mean? (\textquotedblleft{}A house\textquotedblright{}.) Say it once more.
\end{tcolorbox}
\section[\texorpdfstring{\ar{غرفة} (ghurfa)}{غرفة (ghurfa)}]{\ar{غرفة} (ghurfa) — a room}

\label{lesson:AR-C69-ghurfa}



\begin{quote}
Before the new one: say the Arabic for chicken, then the Arabic for a house.
\end{quote}

\subsection*{You'll want to know: \ar{غرفة}}
\textbf{\ar{غرفة}} — \emph{ghurfa} — \textquotedblleft{}a room\textquotedblright{}.

\begin{cousinweb}
\textbf{\ar{غرفة}} (\emph{ghurfa}) is built on the native Arabic root \emph{gh-r-f}, \textquotedblleft{}to ladle out water with the hand\textquotedblright{} (Lane). Lane gives \emph{ghurfa} as a handful of water, and also as an upper chamber, the sense that became \textquotedblleft{}a room\textquotedblright{}.
\end{cousinweb}

\begin{grammarlens}[title={What you've built}]
One more word for this chapter.
\end{grammarlens}

\subsection*{Your turn}
\begin{itemize}
\raggedright
  \item \emph{Say it:} \emph{ghurfa}
  \item \emph{Say it:} \emph{ghurfa}, once more
  \item \emph{Say it:} say \emph{ghurfa}
  \item \emph{Recall:} say the Arabic for chicken, then the Arabic for a house, then say \emph{ghurfa} again
\end{itemize}

\begin{tcolorbox}[breakable,colback=teal!4,colframe=teal!35!black,title={Before you move on}]
What does \ar{غرفة} mean? (\textquotedblleft{}A room\textquotedblright{}.) Say it once more.
\end{tcolorbox}
\section[\texorpdfstring{\ar{نافذة} (nāfidha)}{نافذة (nāfidha)}]{\ar{نافذة} (nāfidha) — a window}

\label{lesson:AR-C69-nafidha}



\begin{quote}
Before the new one: say the Arabic for a house, then the Arabic for a room.
\end{quote}

\subsection*{You'll want to know: \ar{نافذة}}
\textbf{\ar{نافذة}} — \emph{nāfidha} — \textquotedblleft{}a window\textquotedblright{}.

\begin{cousinweb}
\textbf{\ar{نافذة}} (\emph{nāfidha}) is built on the native Arabic root \emph{n-f-dh}, \textquotedblleft{}to pass through, to pierce\textquotedblright{} (Lane). \emph{Nāfidha} is the feminine active participle, \textquotedblleft{}one that passes through\textquotedblright{}; the window sense is not among Lane's.
\end{cousinweb}

\begin{grammarlens}[title={What you've built}]
One more word for this chapter.
\end{grammarlens}

\subsection*{Your turn}
\begin{itemize}
\raggedright
  \item \emph{Say it:} \emph{nāfidha}
  \item \emph{Say it:} \emph{nāfidha}, once more
  \item \emph{Say it:} say \emph{nāfidha}
  \item \emph{Recall:} say the Arabic for a house, then the Arabic for a room, then say \emph{nāfidha} again
\end{itemize}

\begin{tcolorbox}[breakable,colback=teal!4,colframe=teal!35!black,title={Before you move on}]
What does \ar{نافذة} mean? (\textquotedblleft{}A window\textquotedblright{}.) Say it once more.
\end{tcolorbox}
